\section{Hintikka: normal forms and constituent-conditioned induction}
\label{sec:annex-p15-hintikka-inductive-logic}

The normal-form arguments use the declared 1953-based conventions.
The probability arguments establish the specified reconstruction of
formulas in Niiniluoto's historical author copy. These local results do
not establish the full K-dimensional necessity/representation theorem.

\begin{annexitem}{annex:p15.hintikka-monadic-normal-form}{thm:p15-hintikka-monadic-normal-form}{Kinds, constituents and monadic elimination}
\begin{proof}
Every assignment of the \(k\) unary predicates to one object selects
exactly one of their \(2^k\) complete sign conjunctions. Distinct
conjunctions differ on one sign, so cannot both hold. Their disjunction
is a Boolean tautology. For any nonempty model define its support
\(S\) to be the set of profiles realized by some object. The model
satisfies exactly \(C_S\). Conversely each nonempty \(S\) has a
model consisting of one representative of each selected profile.
Thus those constituents are consistent, exhaustive and disjoint.

Here is a quantifier-elimination proof of the normal-form assertion,
including arbitrarily nested monadic formulas. Put
\(E_i=\exists x Q_i(x)\). By induction on formula construction,
every formula with free variables \(y_1,\ldots,y_r\) is a Boolean
combination of their profile assertions \(Q_i(y_j)\) and the closed
sentences \(E_i\). An atomic predicate is the disjunction of the
profiles in which its sign is positive; Boolean steps preserve the
claim. For the existential step, suppose the inductive expression for
\(\varphi(x,\boldsymbol y)\) is known. On fixing a profile \(i\)
for \(x\), replace its profile assertions by their zero/one truth
values, leaving a Boolean expression \(\varphi_i(\boldsymbol y,E)\).
Then
\[
 \exists x\varphi(x,\boldsymbol y)
 \ \equiv\ \bigvee_{i=1}^K
       [E_i\land\varphi_i(\boldsymbol y,E)].
\]
An actual witness has a unique profile and proves the forward direction;
\(E_i\) supplies a witness for the reverse direction. No equality
relation connects \(x\) to another variable, so no further information
about that witness is needed. Universal quantification is handled as
\(\neg\exists x\neg\). This finishes the induction.

A closed sentence is consequently a Boolean combination of the \(E_i\).
Expand that Boolean function into its complete sign conjunctions. The
all-negative conjunction is false on the nonempty domain. Every other
conjunction is exactly \(C_S\), since exclusion of the remaining
profiles is equivalent to the displayed universal restriction. This
gives existence of the normal form. Disjointness and consistency give
uniqueness: omitting or adding one constituent changes truth in its
representative model.

The two displayed quantifier tests follow from the profile partition:
all realized profiles must be allowed for the universal, and at least
one must be allowed for the existential. A sample word realizes every
profile in \(T\), so compatibility requires \(T\subseteq S\).
Conversely that containment permits extending the sample by witnesses
for the remaining selected profiles. Choose any subset of the
\(K-c\) unobserved profiles to obtain the \(2^{K-c}\) compatible
constituents, choosing exactly \(w-c\) for the binomial count.
\end{proof}
\end{annexitem}

\begin{annexitem}{annex:p15.hintikka-information-laws}{thm:p15-hintikka-information-laws}{Information identities and a separating corroboration countercase}
\begin{proof}
For positive \(P(h\land e)\) and \(P(e)\),
\[
 -\log P(h\land e)+\log P(e)=-\log P(h\mid e).
\]
Subtracting conditional surprise from prior surprise gives
\(\log[P(h\mid e)/P(h)]\), positive exactly when its ratio
exceeds one. For content, the additive difference is
\(P(e)-P(h\land e)=P(e)[1-P(h\mid e)]\).
Subtracting it from prior content gives
\(1-P(h)-P(e)+P(h\land e)=1-P(h\lor e)\).
Conditional transmitted content instead gives
\([1-P(h)]-[1-P(h\mid e)]=P(h\mid e)-P(h)\).
These identities retain the different additive and conditional content
definitions rather than identify them by their shared name.

Global entailment means inclusion of the constituent index sets:
\(I_{h_1}\subseteq I_{h_2}\). A sum of nonnegative posterior
masses over the first set is at most the sum over the second. A minimum
over the smaller nonempty set is at least the minimum over the larger.
This proves the two opposite orderings. Under concentration, choose a
stage after which \(P(C_T\mid e)>1/2\); the sum of every other
constituent's posterior is below \(1/2\), so each other individual
posterior is below \(1/2\). A generalization different from \(C_T\)
either omits \(T\) or includes some other constituent. In the first
case its minimum is below \(1/2\); in the second its minimum is
bounded by that other constituent's probability, also below \(1/2\).
Thus \(C_T\) maximizes corroboration.

For the local-entailment failure choose two different constituents
\(C_1,C_2\), take evidence \(e=C_1\), and let
\(h_1=C_1\lor C_2\), \(h_2=C_1\). Positive priors make the
conditioning defined. Relative to \(e\), both hypotheses are true,
so \(e\models h_1\to h_2\). Nevertheless the literal definition
gives \(\operatorname{corr}(h_1,e)=\min(1,0)=0\) and
\(\operatorname{corr}(h_2,e)=1\). This is the claimed countercase
to transferring the global inequality to the evidence-relative claim.
\end{proof}
\end{annexitem}

\begin{annexitem}{annex:p15.hintikka-cognitive-utility}{thm:p15-hintikka-cognitive-utility}{What the specified payoffs actually imply}
\begin{proof}
Write \(p=P(g\mid e)\) and \(r=P(g)\). The indicator of
\(g\)'s truth is one exactly on the constituent answers included
in its normal form, so its expectation is their posterior mass sum
\(p\). A tautology has value one and no truth indicator can exceed
one, proving the first maximization statement.

The second payoff has expectation
\(p(1-r)+(1-p)(-r)=p-r\). The weighted false payoff instead has
expectation
\[
 p(1-r)+(1-p)(-qr)
   =p[1-(1-q)r]-qr.
\]
Its difference from \(p-qr\) is \(-(1-q)pr\), so equality
holds when \(q=1\), or in the degenerate cases \(p=0\) or
\(r=0\), but not generally. With \(r=1/4,p=3/4,q=1/2\)
the true-payoff contribution is \(9/16\), the false-payoff
contribution is \(-1/32\), giving \(17/32\). The other expression
is \(3/4-1/8=5/8\). If the true payoff were \(1-qr\) instead,
its expectation would be \(p(1-qr)-(1-p)qr=p-qr\); that
calculation expressly uses a changed payoff definition.
\end{proof}
\end{annexitem}

\begin{annexitem}{annex:p15.hintikka-truthlikeness-limit}{thm:p15-hintikka-truthlikeness-limit}{Distance, expectation and literal versus approximate truth}
\begin{proof}
For each target \(h_i\), at least one different answer exists and
has positive distance, so \(\sum_j\Delta_{ij}>0\). All terms
in both distances are nonnegative. With positive weights, truthlikeness
equals one exactly when both distances are zero. Sum distance zero
forces every answer in \(I_g\) to have zero distance from \(h_i\),
so \(I_g=\{i\}\) since it is nonempty and distances vanish only
on the diagonal. Conversely that singleton makes both distances zero.

The finite expected-value formula has only finitely many fixed summands.
If the posterior at \(h_*\) tends to one and the others tend to
zero, its limit is the single value \(\operatorname{Tr}(g,h_*)\).
Every answer in \(I_g\) has minimum distance zero from \(g\),
so it belongs to every nonnegative-tolerance blurred answer. Summing
posterior masses proves \(\operatorname{PAT}\geq P(g\mid e)\).
If an answer is outside \(I_g\), its distance to every one of the
finitely many members is positive, and their minimum is positive.
There are only finitely many such outside answers. Below their smallest
positive minimum distance the blurred set is exactly \(I_g\), giving
the asserted zero-tolerance limit (indeed eventual equality).

For the stated two-answer example \(I_g=\{1\}\) and target 2,
the minimum distance is \(1/4\), while the sum-distance numerator
and denominator are both \(1/4\), giving sum distance one.
Thus \(\operatorname{Tr}(g,h_2)=1-(1/4)(1/4)-(1/4)=11/16\).
If evidence makes \(h_2\) certain, this is also expected truthlikeness,
but \(P(g\mid e)=0\). At tolerance \(1/4\), the blurred answer
includes \(h_2\) and therefore has probability one. This proves the
worked separation rather than relying on terminology alone.
\end{proof}
\end{annexitem}

\begin{annexitem}{annex:p15.hintikka-depth-normal-form}{thm:p15-hintikka-depth-normal-form}{The recursive bounded-degree construction}
\begin{proof}
All lists in the construction are finite: at level \(d\) there are
finitely many tuples over finitely many variables and finitely many
predicate symbols; the recursive step forms a finite Boolean list from
a finite collection. For any structure and assignment to the first
\(a\) variables, each listed component has a definite truth value.
Exactly one complete sign conjunction matches those values. Distinct
conjunctions disagree on a component. Hence the descriptions partition
the semantic assignments, including the possibility that some cells of
this partition are empty.

We prove normalizability by descending induction on \(a\), with an
inner induction on Boolean construction. At \(a=d\), all formulas
in question are quantifier-free, so the ordinary finite truth-table
expansion over \(A_d\) gives their normal forms. Suppose the claim
holds at \(a+1\). Atomic formulas at level \(a\) occur in
\(A_a\), and Boolean combinations can be expanded into complete
conjunctions over the level-\(a\) component list. For a quantified
formula, rename its bound variable to \(x_{a+1}\), avoiding capture;
its matrix has free variables among the first \(a+1\) and nesting
at most \(d-(a+1)\). By induction the matrix is equivalent to a
disjunction \(\bigvee_{\tau\in J}\tau\) over
\(\mathcal T_{a+1}\). Existential quantification distributes over
this finite disjunction, producing
\(\bigvee_{\tau\in J}\exists x_{a+1}\tau\).
Each of these existential components is explicitly in the level-\(a\)
list. Universal quantification reduces to negated existential
quantification. The inner Boolean induction now completes the step.
This proof gives an algorithm: recursively expand each matrix and
then apply a finite Boolean truth-table expansion.

The algorithm enumerates component truth values syntactically; it never
asserts that an arbitrary combination is realizable. For example the
description containing both \(\forall xR(x,x)\) and
\(\exists x\neg R(x,x)\), when expanded in the construction's
components, is inconsistent. Existence of an expansion therefore does
not provide a satisfiability filter or probability weights on its
consistent part. No unrestricted relational probability construction
is deduced by this theorem.
\end{proof}
\end{annexitem}

\begin{annexitem}{annex:p15.hintikka-countable-realization}{thm:p15-hintikka-countable-realization}{An explicit infinite realization, not just finite tables}
\begin{proof}
For fixed \(S\), the empty-word likelihood is one. Dividing the
likelihood of a word extended by cell \(i\in S\) by that word's
likelihood gives \((n_i+\lambda_w/w)/(n+\lambda_w)\), since only
one rising factorial acquires a factor. These ratios are positive and
sum to
\((\sum_{i\in S}n_i+w\lambda_w/w)/(n+\lambda_w)=1\).
Cells outside \(S\) have zero likelihood. Induction on word length
therefore proves normalization and deletion consistency. The likelihood
depends only on the counts, proving exchangeability.

For explicit existence use ordinary Lebesgue probability on \([0,1)\).
Partition this interval into finitely many half-open intervals of lengths
\(\pi_S\), one per nonempty support. Within a support interval,
partition recursively into subintervals in increasing cell order, with
lengths equal to the parent length times the displayed conditional
ratios. The children exhaust the parent. Thus each real point chooses
a support and a unique infinite word; half-open conventions settle
endpoints. Equivalently one may discard the countable set of subdivision
endpoints, a null set, without affecting any probability. The inverse
image of a finite word with selected support \(S\) is an interval
of length \(\pi_SL_S(e)\), so the coordinate maps are measurable
and have precisely the desired cylinder probabilities. Boolean formulas
are measurable, and quantifiers over the named domain \(\mathbb N\)
give countable unions or intersections of such events. This supplies an
actual probability interpretation for the monadic sentences, not merely
a putative extension of a finite array.

Fix a selected \(S\) and one \(i\in S\), and put
\(a=\lambda_w/w\). The conditional probability of avoiding cell
\(i\) through \(N\) observations is
\[
 \prod_{t=0}^{N-1}\left(1-\frac a{t+\lambda_w}\right).
\]
For \(w=1\), its first factor is zero. For \(w>1\), all factors
are positive and \(1-u\leq e^{-u}\) bounds this product by
\(\exp[-a\sum_{t<N}(t+\lambda_w)^{-1}]\).
The sum diverges: for sufficiently large dyadic blocks
\(2^b\leq t<2^{b+1}\), their \(2^b\) terms are each at least
\(1/(2^{b+1}+\lambda_w)\), so each block contributes at least
\(1/3\). The product therefore tends to zero. Continuity of
probability on decreasing events shows that cell \(i\) is eventually
witnessed with probability one. A finite union of the zero-probability
failure events for \(i\in S\) still has probability zero. No cell
outside \(S\) occurs, by the recursive construction. Hence the
realized word satisfies exactly \(C_S\) almost surely conditional on
the support choice. Its unconditional mass is exactly \(\pi_S\).

By the monadic normal form, each generalization is a finite disjoint
union of these constituent events. Its probability is their mass sum,
positive whenever the union is nonempty. In particular a consistent
proper universal exclusion retains at least one nonempty support and
has positive probability. This proves the infinite assertion as well
as the finite likelihood formulas.
\end{proof}
\end{annexitem}

\begin{annexitem}{annex:p15.hintikka-combined-calculation}{thm:p15-hintikka-combined-calculation}{Gamma factors, support counts, predictions and worked cases}
\begin{proof}
For completeness, the positive-real gamma integral
\(\Gamma(u)=\int_0^\infty t^{u-1}e^{-t}\,dt\) exists for
\(u>0\). Near zero its integrand is bounded by \(t^{u-1}\),
which is integrable. At infinity choose an integer \(M\geq u-1\);
the exponential series gives \(t^M\leq2^MM!e^{t/2}\), bounding
the tail by a constant times \(e^{-t/2}\). The integral is strictly
positive. Integration by parts, with boundary terms vanishing by the
same exponential bound and \(t^u\to0\) at zero, gives
\(\Gamma(u+1)=u\Gamma(u)\). Also \(\Gamma(1)=1\), whence
\(\Gamma(r+1)=r!\) for integers \(r\geq0\), and
\((u)_r=\Gamma(u+r)/\Gamma(u)\). Thus the prior formulas are
positive finite quantities and their finite normalizing sums are nonzero.

For \(\lambda_w=w\), each numerator rising factorial is
\((1)_{n_i}=n_i!\), and \((w)_n=(n+w-1)!/(w-1)!\).
This proves the likelihood formula. Multiplying by the prior weight
\(\Gamma(\alpha+w)/\Gamma(w)\) cancels \(\Gamma(w)\),
leaving \(a_w(n,\alpha)\prod_{i\in T}n_i!\), apart from the
common prior normalizer. Bayes normalization cancels both common
factors. The count of compatible supports of width \(w\) is
\(\binom{K-c}{w-c}\), proving the denominator and both width
interpretations. An exclusion law retains only supports contained in
the \(K-b\) allowed cells, giving the stated numerator. If an
excluded cell has been observed, every compatible support violates the
law, so its posterior is zero.

For an observed target cell, every compatible support contains it, so
average \((n_i+1)/(n+w)\) using the width multiplicities. For an
unobserved target, only supports which contain that target contribute;
after choosing it there remain \(K-c-1\) available cells, from which
to choose \(w-c-1\). Average their \(1/(n+w)\) predictions.
This establishes the distinct unobserved-cell formula.

For the first four-cell example the common denominator 2520 gives
unnormalized weights 21,7,7,3, whose sum is 38. Thus the observed-cell
prediction is
\(3(28+16+3)/304=141/304\). For each unobserved cell it is
\(8/304+3/304=11/304\); the two observed and two unobserved
probabilities sum to one. For the second example widths 3 and 4 have
weights \(1/360,1/840\), hence relative weights 7 and 3.
The target prediction is \(3(1/10+3/80)=33/80\), distinct from
\(141/304\). These examples preserve the same \(n,n_1\), while
changing \(c\), as claimed.
\end{proof}
\end{annexitem}

\begin{annexitem}{annex:p15.hintikka-generalization-convergence}{thm:p15-hintikka-generalization-convergence}{Concentration, arbitrary observed frequencies and finite bounds}
\begin{proof}
Start with the constant-\(\lambda\) family. For any compatible
\(S\supsetneq T\) of width \(w>c\), put
\(a=\lambda/w\), \(b=\lambda/c\), so \(0<a<b\).
Cancellation of the common denominator gives
\[
 \frac{L_S(e)}{L_T(e)}
 =\prod_{i\in T}\prod_{j=0}^{n_i-1}\frac{a+j}{b+j}.
\]
Each factor is at most one. Some count is at least \(n/c\), so
the ratio is bounded above by
\[
 \prod_{j=0}^{\lfloor n/c\rfloor-1}
       \left(1-\frac{b-a}{b+j}\right)
 \leq\exp\left[-(b-a)\sum_{j<\lfloor n/c\rfloor}\frac1{b+j}\right]
 \longrightarrow0.
\]
The harmonic divergence follows by the dyadic-block argument already
given. Multiplying by the fixed positive prior ratio does not change
the limit. There are finitely many compatible supports. Dividing all
posterior weights by \(\pi_TL_T(e)\) shows that the denominator
is one plus a finite sum tending to zero, proving concentration.

In the combined family gamma recurrence gives exactly
\[
 a_w/a_c=\prod_{j=c}^{w-1}(\alpha+j)/(n+j).
\]
Each product tends to zero for fixed finite \(\alpha\). The same
finite-sum normalization proves concentration. A generalization's
probability is a sum over its constituent normal form, so its limit
is one or zero according as that form contains \(C_T\) or not.
At any finite stage every compatible constituent has a positive
likelihood and a positive prior, so two surviving supports still have
strictly positive posterior probability. Logical entailment has not
been obtained by passing to a high finite posterior.

For the bound, each factor in \(a_w/a_c\) is at most
\(r=(\alpha+K-1)/(n+c)\). Therefore
\[
 D_{n,c}/a_c\leq
 \sum_{j=0}^{K-c}\binom{K-c}j r^j=(1+r)^{K-c}.
\]
Taking reciprocals gives the claimed posterior bound. For \(m=K-c>0\),
\((1+r)^{-m}\geq1-\varepsilon\) is equivalent to
\(r\leq(1-\varepsilon)^{-1/m}-1\), giving the sufficient
sample-size inequality. If \(m=0\), only \(T\) survives.
\end{proof}
\end{annexitem}

\begin{annexitem}{annex:p15.hintikka-prediction-relevance}{thm:p15-hintikka-prediction-relevance}{Frequency matching and strictly positive instantial relevance}
\begin{proof}
For constant \(\lambda\), on a support containing the target,
\[
 \frac{n_i+\lambda/w}{n+\lambda}-\frac{n_i}n
 =\frac{\lambda(n/w-n_i)}{n(n+\lambda)}.
\]
Since \(0\leq n_i\leq n\), its absolute value is at most
\(\lambda/n\). A support not containing the target has
\(n_i=0\) and conditional prediction zero; the same bound holds.
Averaging preserves the bound. For the combined family replace the
numerator by \(n-wn_i\); its absolute value is at most
\((K+1)n\), giving the displayed, conservative bound after division.

For relevance write \(v_S=P(C_S\mid e)\) and let \(q_S\)
be the target's conditional next probability, zero if the support
excludes it. The unconditional next probability
\(p=\sum_Sv_Sq_S\) is positive because the full support has positive
weight. After a positive instance, Bayes's formula gives support
weights \(v_Sq_S/p\). On a support containing the target its
updated conditional probability is \(q_S+\delta_S\), where
\[
 \delta_S=\frac{n+\lambda_w-n_i-\lambda_w/w}
                 {(n+\lambda_w)(n+1+\lambda_w)}\geq0.
\]
The resulting unconditional change is
\[
 p'-p=\frac{\sum_Sv_Sq_S^2-p^2}{p}
       +\frac{\sum_Sv_Sq_S\delta_S}{p}.
\]
The first numerator is the finite weighted variance of the \(q_S\),
and is nonnegative by expansion of
\(\sum_Sv_S(q_S-p)^2\). The second is nonnegative. For the full
support, \(K>1\) and \(\lambda_K>0\) make
\(n-n_i+\lambda_K(1-1/K)>0\); all its other factors are positive.
Thus the second sum is strictly positive, proving the strict increase.
This proof uses the declared conditional family, not an unproved
implication from a general convergence axiom.
\end{proof}
\end{annexitem}

\begin{annexitem}{annex:p15.hintikka-caution-limits}{thm:p15-hintikka-caution-limits}{Parameter endpoints and two different orders of limiting}
\begin{proof}
For \(w>c\), gamma recurrence gives the relative prior ratio
\[
 r_w/r_c=\frac{\Gamma(c)}{\Gamma(w)}
                \prod_{j=c}^{w-1}(\alpha+j).
\]
It is strictly increasing in \(\alpha\geq0\). For fixed evidence,
the analogous posterior ratio \(a_K/a_w\) is a positive constant
times \(\prod_{j=w}^{K-1}(\alpha+j)\), hence tends to infinity
for every \(w<K\). Finite normalization concentrates on full width.
The mixture predictive rule therefore tends to its full-width value
\((n_i+1)/(n+K)\). On the other hand the preceding concentration
theorem gives limit one for \(C_T\) for every fixed finite
\(\alpha\). If \(c<K\), concentration on full width gives
limit zero for \(C_T\) for every fixed finite sample. Applying
the remaining outer limits yields exactly the two displayed values.

Within a fixed support, dividing numerator and denominator of
\((n_i+\lambda/w)/(n+\lambda)\) by \(\lambda\) gives limit
\(1/w\) as \(\lambda\to\infty\); letting \(\lambda\downarrow0\)
with \(n>0\) gives \(n_i/n\). A finite ordered word's likelihood
is the product of its finitely many step probabilities, so the first
limit gives \(w^{-n}\). With equal nonempty-support priors, this
is the Jerusalem realization. The proof asserts this declared limit,
not literal equivalence of all finite-domain allocations in unheld
original sources.
\end{proof}
\end{annexitem}

\begin{annexitem}{annex:p15.hintikka-two-cell-comparison}{thm:p15-hintikka-two-cell-comparison}{The law probability and next-instance probability are different}
\begin{proof}
The three support priors are \(1/3\). Evidence containing only
\(A\) eliminates the \(\neg A\)-only support. Conditional on
the \(A\)-only support the likelihood is one. Conditional on full
support the combined likelihood is \(n!/(n+1)!=1/(n+1)\).
The posterior of the \(A\)-only constituent, which is exactly
\(H\) on the nonempty domain, is
\(1/[1+1/(n+1)]=(n+1)/(n+2)\). The full-support posterior is
\(1/(n+2)\). Only that latter support can predict a contrary
instance, with conditional probability \(1/(n+2)\); the unconditional
contrary-instance probability is its square. This proves the combined
prediction.

For uniform independent sampling on full support its likelihood is
\(2^{-n}\). Normalizing against the \(A\)-only likelihood one
gives law probability \(2^n/(2^n+1)\); the next contrary probability
is \(1/[2(2^n+1)]\), giving the other formula. Substitute \(n=2\)
to obtain all four numerical values.

For Carnap's equal-structure measure on \(N\) individuals and two
profiles, the all-\(A\) state has probability \(1/(N+1)\),
while the report of \(n\) specified positive individuals has
probability \(1/(n+1)\). Hence the finite universal probability is
\((n+1)/(N+1)\), tending to zero as \(N\to\infty\) with the
sample fixed. His next-instance predictive value is
\((n+1)/(n+2)\), by the equal-structure predictive calculation.
Finally the sample empirical frequency is \(n/n=1\). That numerical
fact establishes the straight guess but provides no prior event mass
for an infinite conjunction, so no Reichenbach law probability is
inferred from it.
\end{proof}
\end{annexitem}

\begin{annexitem}{annex:p15.hintikka-evidential-success}{thm:p15-hintikka-evidential-success}{An explicit selection failure and the added sampling assumption}
\begin{proof}
Take a countable domain containing infinitely many \(A\) objects and
at least one \(\neg A\) object. Enumerate only the \(A\) objects
as evidence. Every report is true, and its observed support is always
\(\{A\}\). The two-cell formulas tend to one for the universal
law, but the domain's contrary object falsifies that law. Thus true
reports plus infinite sample size do not imply truth convergence.

If all reports are true and each truly instantiated cell is eventually
witnessed, observed support cannot exceed true support and eventually
contains it. Because there are finitely many cells, a finite maximum
of their first witnessing times gives eventual equality. The
concentration theorem now applies to the true constituent.

Under independent draws with a fixed positive probability \(p_i\)
for each true cell, the probability of missing cell \(i\) through
\(N\) draws is \((1-p_i)^N\), tending to zero. Continuity on
decreasing events shows that its probability of never appearing is zero.
The finite union over all true cells still has probability zero, proving
almost-sure witnessing. Nonetheless a stream selecting only one cell
is logically possible when another exists; it is just null under these
sampling assumptions. Every finite one-cell prefix also has a
continuation which eventually observes another cell. No such prefix
certifies exhaustion of the actual world's support.
\end{proof}
\end{annexitem}
